\documentclass[11pt,a4paper]{article}
% =====================================================================
%  Shared preamble for the Part V lecture notes (lectures 19-22).
%
%  These four lectures sit outside Geron, so the notes ARE the primary
%  source and are examined on the same terms as the textbook parts.
%  Everything here is chosen to make that readable on paper: one column,
%  generous measure, and the same six-colour palette as the slides so a
%  student moving between the two is not re-learning what red means.
%
%  Build with pdflatex (twice, for the toc and the refs):
%      cd notes && latexmk -pdf lecture-19.tex
%  or  make -C notes
% =====================================================================
\usepackage[T1]{fontenc}
\usepackage[utf8]{inputenc}
\usepackage{newtxtext,newtxmath}      % Times-like: prints well, reads well
\usepackage[scaled=0.86]{inconsolata} % code, at a size that sits beside Times
\usepackage{microtype}
% newtx defines \Bbbk, \openbox, \proof and \endproof; amssymb and amsthm
% then redefine them and abort. Freeing the four names before those two
% packages load is the standard fix, and it costs nothing here: we use
% amsthm only for \newtheorem.
\let\Bbbk\relax
\usepackage{amsmath,amssymb}
\let\openbox\relax
\let\proof\relax
\let\endproof\relax
\usepackage{amsthm}
\usepackage{booktabs}
\usepackage{enumitem}
\usepackage{xcolor}
\usepackage{tcolorbox}
\usepackage{titlesec}
\usepackage{graphicx}
\usepackage[hidelinks]{hyperref}
\tcbuselibrary{skins,breakable}

% ---- the course palette, from assets/css/custom.css -------------------
\definecolor{aimlprimary}{HTML}{0B3D62}   % deep blue  - structure
\definecolor{aimlaccent} {HTML}{C0392B}   % brick red  - failures
\definecolor{aimlsuccess}{HTML}{14663A}   % green      - repairs
\definecolor{aimlmath}   {HTML}{6C3483}   % purple     - the derivation
\definecolor{aimlmuted}  {HTML}{4B5563}
\definecolor{aimlrule}   {HTML}{B0BCC7}
\definecolor{aimlsoft}   {HTML}{F4F7F9}

\usepackage[a4paper,margin=32mm,top=30mm,bottom=30mm]{geometry}
\setlength{\parskip}{0.45\baselineskip}
\setlength{\parindent}{0pt}
\linespread{1.06}

% ---- headings --------------------------------------------------------
\titleformat{\section}{\Large\bfseries\color{aimlprimary}}{\thesection}{0.7em}{}
\titleformat{\subsection}{\large\bfseries\color{aimlprimary}}{\thesubsection}{0.7em}{}
\titleformat{\subsubsection}{\bfseries\color{aimlmuted}}{}{0em}{}
\titlespacing*{\section}{0pt}{1.6\baselineskip}{0.5\baselineskip}
\titlespacing*{\subsection}{0pt}{1.2\baselineskip}{0.35\baselineskip}

% ---- boxes -----------------------------------------------------------
% One box per role, and the role is the colour. A student who has seen the
% slides already knows what each one means before reading a word of it.
\newtcolorbox{keybox}[1][]{
  breakable, enhanced, colback=aimlsoft, colframe=aimlprimary,
  boxrule=0.4pt, left=8pt, right=8pt, top=6pt, bottom=6pt,
  fonttitle=\bfseries\small, coltitle=white, colbacktitle=aimlprimary,
  attach boxed title to top left={yshift=-2mm, xshift=4mm},
  boxed title style={boxrule=0pt, arc=1pt}, #1}

\newtcolorbox{failbox}[1][]{
  breakable, enhanced, colback=white, colframe=aimlaccent,
  boxrule=0.9pt, left=8pt, right=8pt, top=6pt, bottom=6pt,
  fonttitle=\bfseries\small, coltitle=white, colbacktitle=aimlaccent,
  attach boxed title to top left={yshift=-2mm, xshift=4mm},
  boxed title style={boxrule=0pt, arc=1pt}, #1}

\newtcolorbox{derivbox}[1][]{
  breakable, enhanced, colback=white, colframe=aimlmath,
  boxrule=0.6pt, left=8pt, right=8pt, top=6pt, bottom=6pt,
  fonttitle=\bfseries\small, coltitle=white, colbacktitle=aimlmath,
  attach boxed title to top left={yshift=-2mm, xshift=4mm},
  boxed title style={boxrule=0pt, arc=1pt}, #1}

% An exercise. Deliberately the same shape as the notebook's `try` field:
% one change, and what should happen to the output.
\newtcolorbox{trybox}[1][]{
  breakable, enhanced, colback=white, colframe=aimlsuccess,
  boxrule=0.5pt, left=8pt, right=8pt, top=6pt, bottom=6pt,
  fonttitle=\bfseries\small, coltitle=white, colbacktitle=aimlsuccess,
  attach boxed title to top left={yshift=-2mm, xshift=4mm},
  boxed title style={boxrule=0pt, arc=1pt}, #1}

% ---- small semantics -------------------------------------------------
\newcommand{\scope}[1]{\textcolor{aimlmuted}{\small #1}}
\newcommand{\term}[1]{\textbf{#1}}
\newcommand{\code}[1]{\texttt{\small #1}}
\newcommand{\lecture}[1]{Lecture~#1}

\newtheorem{definition}{Definition}[section]
\newtheorem{proposition}[definition]{Proposition}

% ---- title block -----------------------------------------------------
% \notestitle{19}{Information retrieval: the lexical foundation}{SciFact (BEIR)}
\newcommand{\notestitle}[3]{%
  \thispagestyle{empty}
  \begin{flushleft}
    {\color{aimlmuted}\small\sffamily
     APPLICATIONS OF MACHINE LEARNING \quad$\cdot$\quad
     BSc Mathematics of Artificial Intelligence}\\[2mm]
    {\color{aimlrule}\rule{\textwidth}{0.8pt}}\\[4mm]
    {\color{aimlmuted}\large Lecture #1 \quad$\cdot$\quad Part V \quad$\cdot$\quad
     Extended notes}\\[3mm]
    {\Huge\bfseries\color{aimlprimary}\baselineskip=1.15\baselineskip #2\par}\vspace{4mm}
    {\color{aimlmuted}\large Dataset: #3}\\[3mm]
    {\color{aimlrule}\rule{\textwidth}{0.8pt}}
  \end{flushleft}
  \vspace{2mm}
  \begin{keybox}[title={Why these notes exist}]
    Lectures 19--22 sit outside G\'eron, the course's primary text. For those
    four lectures \emph{these notes are the primary source}, and they are
    examined on exactly the same terms as the textbook parts. Everything in
    the slide deck for this lecture is here, with the arguments written out at
    length rather than compressed onto a slide; the numbers are the ones the
    accompanying notebook prints, and every one of them is a measurement on
    the corpus named above rather than a figure quoted from a paper.
  \end{keybox}
  \vspace{1mm}
  {\small\tableofcontents}
  \vspace{2mm}
  \noindent{\color{aimlrule}\rule{\textwidth}{0.4pt}}
  \vspace{1mm}
}


\begin{document}
\notestitle{22}{Recommender systems: neural, and evaluated honestly}{MovieLens 1M, binarised}

\section{From ratings to rankings}

\lecture{21} asked what a user would rate a film. This lecture asks what to
show them. They are not the same question, and the second is the one a
recommender is actually deployed to answer.

\begin{center}
\small
\begin{tabular}{p{0.14\textwidth}p{0.33\textwidth}p{0.4\textwidth}}
\toprule
 & \textbf{\lecture{21}} & \textbf{\lecture{22}} \\
\midrule
Question      & what would they rate it? & what should we show? \\
Output        & a number per pair & an ordered list of ten \\
Evaluated on  & ratings they gave & items they had not yet touched \\
Metric        & RMSE & HR@10, NDCG@10 --- \lecture{19}'s \\
Data          & explicit ratings & implicit interactions \\
\bottomrule
\end{tabular}
\end{center}

Row three changes everything. The evaluation set is no longer a sample of what
the user did; it is a claim about what they \emph{would} do, on items they
never saw.

\subsection{Precisely why RMSE fails here}

\lecture{21}'s RMSE of 0.8587 was computed on held-out ratings --- pairs where
the user chose to watch the film \emph{and} chose to rate it.

\begin{itemize}[leftmargin=*,nosep]
  \item A recommender operates on the complement of that set: films the user
        has not chosen.
  \item RMSE weights every prediction equally, and a recommender shows ten
        items. Being accurate on the three-star films costs nothing and buys
        nothing.
  \item And it is a \emph{pointwise} error, whereas showing ten items is a
        decision about \emph{order}.
\end{itemize}

\begin{keybox}[title={A concrete case}]
A model that predicts every rating 0.3 stars too low has a poor RMSE and a
\emph{perfect} ranking: a constant offset changes no order at all. The two
metrics can disagree completely.
\end{keybox}

\subsection{The data, binarised}

Call a rating of 4 or 5 an interaction and discard the rest:

\begin{center}
\small
\begin{tabular}{lr}
\toprule
Ratings                    & 1{,}000{,}209 \\
Positive interactions      & 575{,}281 \\
As a fraction of ratings   & 57.5\% \\
Cells with no interaction  & 97.4\% \\
\bottomrule
\end{tabular}
\end{center}

Note what the threshold did: a film a user rated 3 is now indistinguishable
from one they never saw. We have thrown away the only genuine negatives in the
dataset, in order to simulate the situation of a real system --- which never
had any.

\subsection{The metrics, from Lecture 19}

\term{HR@10} is the hit rate: is the held-out item in the top ten? For one
held-out item this is recall@10, and it is 0 or 1 per user. \term{NDCG@10} is
the same, discounted by position: $1/\log_2(\text{rank}+1)$ when hit, 0
otherwise. Both come straight from \lecture{19}'s derivation with $|R| = 1$.
Nothing new is needed, which is the point of having done retrieval first.

\begin{derivbox}[title={HR@10, defined precisely}]
\begin{enumerate}[leftmargin=*,nosep]
  \item Hold out exactly one interaction per user; call its item $g$.
  \item Score the candidate set for that user; \emph{remove every item they
        interacted with in the training data}.
  \item Take the top ten. $\text{HR} = 1$ if $g$ is among them, else 0.
  \item Average over users --- macro, one vote each.
\end{enumerate}
Step 2 is where the candidate set is defined, and the candidate set is half of
what this lecture measures.
\end{derivbox}

\subsubsection{Why already-seen items must be excluded}

A small implementation detail with a large effect, and a good example of how a
protocol leaks. Leave them in and the popularity baseline recommends the
blockbusters every heavy user has already watched --- scoring nothing, but
occupying the ten slots. Leave them in for the factorisation and it recommends
what the user already interacted with, because that is what it was trained to
score highly. Which of the two is hurt more depends on the model, so the choice
is \emph{not neutral between them}. We exclude them everywhere, and say so.

\subsection{The baseline that must be in every table}

Recommend the ten most popular items. To everybody. Every time. No
personalisation, no model, no parameters --- minus the items this user has
already interacted with, which is the only thing in it that varies by user.

\begin{keybox}[title={It is not a straw man}]
In the reproducibility literature this baseline beats a substantial fraction of
published neural recommenders, and by the end of this lecture you will see
exactly how that happens. A recommender that does not beat popularity has not
been shown to personalise anything.
\end{keybox}

\section{Two towers}

\subsection{It is Lecture 20, with the query replaced}

\[ \text{score}(u,i) = E_{\text{user}}(u)\cdot E_{\text{item}}(i) \]

\begin{center}
\small
\begin{tabular}{p{0.19\textwidth}p{0.3\textwidth}p{0.36\textwidth}}
\toprule
 & \textbf{\lecture{20}} & \textbf{Here} \\
\midrule
Left tower     & the query text & the user --- id, history, context \\
Right tower    & the document   & the item --- id, genre, text \\
Precomputed    & document vectors & item vectors \\
At request time & encode the query & encode the user \\
\bottomrule
\end{tabular}
\end{center}

Identical architecture, identical indexability argument, identical
approximate-nearest-neighbour problem. \emph{Retrieval and recommendation
converge here}, and that convergence is why Part~V is one part rather than two.

\subsection{What the towers can hold}

Matrix factorisation is the special case where each tower is a \emph{lookup
table}: $E_{\text{user}}(u) = p_u$ and $E_{\text{item}}(i) = q_i$, and there
the resemblance ends. A tower can instead be a network over \emph{features}:
the item's genre, its text, its age; the user's recent history, device, time of
day. Which is what buys the cold start --- a brand-new item has no $q_i$ to
look up, but it does have a genre and a description.

\begin{keybox}[title={The honest summary of ``neural recommenders''}]
The architecture is a dot product of two learned vectors, exactly as before.
What is new is \emph{what the vectors are computed from} --- and on datasets
with no features beyond ids, there is nothing for the network to compute from.
That is a large part of why the reproducibility results came out as they did.
\end{keybox}

\begin{center}
\small
\begin{tabular}{p{0.18\textwidth}p{0.32\textwidth}p{0.34\textwidth}}
\toprule
 & \textbf{Lookup table (MF)} & \textbf{Feature tower} \\
\midrule
New item  & no vector at all & encode its genre and text \\
New user  & no vector at all & encode context and first actions \\
Tail item & estimated from very few interactions & shares parameters with
similar items \\
\bottomrule
\end{tabular}
\end{center}

17.9\% of films here have fewer than twenty interactions, so row three is most
of the catalogue. \emph{This} --- not accuracy on the head --- is the honest
case for a neural recommender, and it is invisible to a metric dominated by
popular items.

\subsection{Training: the softmax over the catalogue}

\[
  \mathcal{L} = -\log
  \frac{\exp(s_{ui})}{\sum_{j\in\text{catalogue}} \exp(s_{uj})}
\]

The natural objective: the observed item should win against every other item.
It is the cross-entropy of \lecture{17} with the catalogue as the classes.

And it is unaffordable. The denominator runs over 3{,}706 films here, and over
tens of millions in a real catalogue --- per interaction, per step. So sample
the denominator. Which introduces a bias, and the bias has a correction, and
the correction is the examinable part.

\begin{derivbox}[title={Sampled softmax and the $\log q$ correction}]
Draw $m$ negatives from a proposal distribution $q$ and sum over those instead.
Popular items are drawn more often, so they are pushed down more often --- the
sample is not the catalogue, and the gradient knows it. Correct by
\[ s'_{uj} = s_{uj} - \log q(j) \]
for each sampled negative. This makes the sampled sum an unbiased estimator of
the full one. Without it, the model over-corrects against popular items and
systematically under-recommends them.

It is the same correction as in word2vec's sampled objectives, and the same
idea as importance weighting: reweight by how likely you were to have drawn the
thing you drew.
\end{derivbox}

\subsubsection{In-batch negatives, once more}

\begin{center}
\small
\begin{tabular}{rrrr}
\toprule
\textbf{Batch $B$} & \textbf{Encoder passes} & \textbf{Scores} &
\textbf{Negatives per user} \\
\midrule
8      & 16       & 64          & 7 \\
32     & 64       & 1{,}024     & 31 \\
128    & 256      & 16{,}384    & 127 \\
512    & 1{,}024  & 262{,}144   & 511 \\
2{,}048 & 4{,}096 & 4{,}194{,}304 & 2{,}047 \\
\bottomrule
\end{tabular}
\end{center}

Exactly \lecture{20}'s table and the same arithmetic --- with one difference
that matters. In-batch negatives are drawn from the \emph{interaction stream},
so an item's sampling probability is proportional to its popularity, which is
precisely the $q(j)$ the correction needs.

\begin{failbox}[title={Skip the correction and popularity bias is in the gradient}]
Not in the data, not in the metric --- in the gradient itself. It is the
cheapest possible bug to introduce and one of the hardest to see afterwards,
because the model looks as though it simply learned to dislike popular items.
\end{failbox}

Where the work really is: not in the architecture, but in the sampling
distribution, the correction, the freshness of the item index --- item vectors
drift as the tower trains, and a stale index is silently wrong, exactly as in
\lecture{20} --- and the protocol you evaluate under. That is the same list as
\lecture{20}'s, and it is why these two lectures are a pair.

\section{The protocol is an experimental variable}

\subsection{Two decisions nobody reports}

To evaluate a recommender you must decide two things, and papers routinely
state neither.

\begin{enumerate}[leftmargin=*]
  \item \term{Which interaction is held out?} One chosen at random from the
        user's history, or their last one in time.
  \item \term{Against what is it ranked?} The whole catalogue, or the held-out
        item against a sample of, say, 100 items the user did not interact
        with.
\end{enumerate}

Two binary choices, four protocols. We are going to run the same model on the
same data under all four.

\begin{trybox}[title={Commit now}]
Write down how much you think the four numbers can differ. A few percent? A
factor of two? Keep the number; you will want it in three pages.
\end{trybox}

\subsubsection{Stating a protocol, in full}

A complete protocol names six things. Most papers name two.

\begin{enumerate}[leftmargin=*,nosep]
  \item How the data was filtered --- here, ratings of 4 or 5.
  \item What was held out --- one interaction per user.
  \item How it was chosen --- at random, or the most recent.
  \item What it is ranked against --- the catalogue, or a sample.
  \item What was excluded from the candidates --- already-seen items.
  \item The metric and cutoff --- HR@10, NDCG@10.
\end{enumerate}

Change any one and the number changes.

\subsection{Decision 1: which interaction is held out}

\begin{center}
\small
\begin{tabular}{p{0.2\textwidth}p{0.32\textwidth}p{0.32\textwidth}}
\toprule
 & \textbf{Random (leave one out)} & \textbf{Temporal (leave last)} \\
\midrule
Held out       & one interaction, chosen uniformly & the user's most recent \\
Training set contains & the user's future & only their past \\
Matches deployment & no & yes \\
In the literature & very common & less common \\
\bottomrule
\end{tabular}
\end{center}

\begin{failbox}[title={The random split is a leak, and it looks like a convention}]
\lecture{2} defined a leak: information in the training data that will not be
available at prediction time. Check the definition. The model is asked to
predict interaction $t$; it was trained on interactions $t{+}1, t{+}2, \ldots$
by that same user; a deployed system has none of them, because they have not
happened.

No column was duplicated, no target was copied, nothing was fitted before the
split. The leak is in the \emph{shape of the evaluation}, which is why it
survived so long as a default.
\end{failbox}

\subsection{Decision 2: against what is it ranked}

\begin{center}
\small
\begin{tabular}{p{0.2\textwidth}p{0.28\textwidth}p{0.36\textwidth}}
\toprule
 & \textbf{Full catalogue} & \textbf{Sampled, 100 negatives} \\
\midrule
Candidates ranked  & 3{,}706 & 101 \\
Cost               & score every item & score 101 \\
Matches deployment & yes --- the system ranks the catalogue & no \\
Why anyone does it & --- & it was expensive, once \\
\bottomrule
\end{tabular}
\end{center}

Sampling was a computational shortcut from a time when scoring a catalogue was
expensive. Ranking one held-out item against 100 random others is a far easier
task than ranking it against 3{,}706 --- and the 100 are drawn uniformly, so
they are mostly obscure films nobody would recommend.

\subsection{The experiment}

One model: a rank-32 factorisation of the binarised interaction matrix --- the
simplest two-tower model there is, with both towers lookup tables. One dataset:
MovieLens 1M, 575{,}281 positive interactions. One baseline: the most popular
items, minus what the user has already seen. Four protocols:
$\{\text{random}, \text{temporal}\} \times \{\text{full}, \text{sampled }100\}$.

Nothing is tuned per protocol, because tuning per protocol is one of the ways
the comparisons in the literature broke. The claim is not that this model is
good; it is that \emph{the protocol moves the number more than the model does},
and that argument is cleanest when the model is one you could have written in
\lecture{21}.

\subsection{The table}

\begin{center}
\small
\begin{tabular}{lrr}
\toprule
\textbf{HR@10} & \textbf{Full catalogue} & \textbf{Sampled 100} \\
\midrule
Temporal $\cdot$ factorisation & 0.0926 & 0.6858 \\
Temporal $\cdot$ popularity    & 0.0404 & 0.4819 \\
Random $\cdot$ factorisation   & 0.2132 & 0.8085 \\
Random $\cdot$ popularity      & 0.0879 & 0.6201 \\
\bottomrule
\end{tabular}
\end{center}

The same model on the same data scores 0.0926 and 0.8085 --- a factor of 8.7.
Nothing about the model changed. Only the two sentences describing how it was
scored.

\begin{center}
\small
\begin{tabular}{lrr}
\toprule
\textbf{NDCG@10} & \textbf{Full catalogue} & \textbf{Sampled 100} \\
\midrule
Temporal $\cdot$ factorisation & 0.0464 & 0.4178 \\
Temporal $\cdot$ popularity    & 0.0196 & 0.2644 \\
Random $\cdot$ factorisation   & 0.1204 & 0.5612 \\
Random $\cdot$ popularity      & 0.0462 & 0.3696 \\
\bottomrule
\end{tabular}
\end{center}

Same shape, a factor of 12.1 across the corners. Worth checking, because ``use
a better metric'' is the natural first response to the HR table --- and it does
not help. The problem is not the metric; it is what the metric was computed
over.

\subsection{Take it apart: the split}

\begin{center}
\small
\begin{tabular}{lrrr}
\toprule
\textbf{HR@10, full catalogue} & \textbf{Temporal} & \textbf{Random} &
\textbf{Ratio} \\
\midrule
Factorisation & 0.0926 & 0.2132 & $\times 2.30$ \\
Popularity    & 0.0404 & 0.0879 & $\times 2.18$ \\
\bottomrule
\end{tabular}
\end{center}

Switching from a temporal to a random split more than doubles both numbers, and
does it to the personalised model and the unpersonalised baseline alike.

\begin{keybox}[title={That last part is the tell}]
If the gain were the model learning something, it would not lift a popularity
list by the same factor. It is \emph{the task getting easier}: predicting a
randomly chosen item from a history you have already seen most of.
\end{keybox}

\subsection{Take it apart: the sampling}

\begin{center}
\small
\begin{tabular}{lrrr}
\toprule
\textbf{HR@10, temporal split} & \textbf{Full} & \textbf{Sampled 100} &
\textbf{Ratio} \\
\midrule
Factorisation & 0.0926 & 0.6858 & $\times 7.4$ \\
Popularity    & 0.0404 & 0.4819 & $\times 11.9$ \\
\bottomrule
\end{tabular}
\end{center}

Sampling is by far the larger effect. HR@10 against 100 random negatives asks
``is the right item in the top ten of \emph{101}?'' --- and ten of 101 is
already a tenth of the candidate set. A random ranker scores 0.0990 under this
protocol and about 0.0027 against the full catalogue: \emph{the floor moved by
a factor of 37}, and every method rides up with it.

\subsubsection{Why sampling flatters unequally}

The 100 negatives are drawn uniformly, and the catalogue is mostly tail: the
top 1\% of films hold 8.7\% of interactions.

\begin{itemize}[leftmargin=*,nosep]
  \item A uniform sample of 100 therefore contains, in expectation, about one
        head film --- the only items a user might plausibly have chosen
        instead.
  \item A popularity model is competing against almost no popular items, so it
        wins easily.
  \item A model whose skill is discriminating \emph{among plausible items} has
        nothing to discriminate.
\end{itemize}

\begin{failbox}[title={So sampled metrics are not merely noisy}]
The distortion depends on what kind of model you have, which is exactly the
condition under which a metric can \emph{reverse an ordering}. A bias that were
uniform would at least preserve rankings; this one does not. That is the
finding the recommender literature had to absorb, and it invalidated a great
many published comparisons.
\end{failbox}

\subsection{How large are modelling gains, for comparison?}

To read those effects you need a sense of what a real improvement looks like in
this field. A new architecture reporting a 5--10\% relative gain on HR@10 is a
publishable result. Our two protocol decisions moved HR@10 by 130\% and 641\%
relative, on an unchanged model.

\begin{keybox}[title={The reproducibility problem, in one comparison}]
The undocumented choices are an order of magnitude larger than the documented
contributions.
\end{keybox}

\subsection{The consequence that ends careers}

\begin{center}
\small
\begin{tabular}{lr}
\toprule
\textbf{System, as reported} & \textbf{HR@10} \\
\midrule
Popularity, random split, sampled 100          & 0.6201 \\
Factorisation, temporal split, full catalogue  & 0.0926 \\
\bottomrule
\end{tabular}
\end{center}

A method with \emph{no personalisation whatsoever}, evaluated generously,
reports 6.7 times the score of a real factorisation evaluated honestly. Both
numbers are correct and both are computed without error. Put them in a table
together, as papers do when they quote a baseline from another paper, and the
comparison is meaningless --- and nothing in either number says so.

\subsection{What is stable across protocols}

\begin{center}
\small
\begin{tabular}{lr}
\toprule
\textbf{Protocol} & \textbf{Factorisation $\div$ popularity} \\
\midrule
Temporal, full catalogue & $\times 2.29$ \\
Random, full catalogue   & $\times 2.42$ \\
Temporal, sampled 100    & $\times 1.42$ \\
Random, sampled 100      & $\times 1.30$ \\
\bottomrule
\end{tabular}
\end{center}

The absolute numbers spanned a factor of 8.7; the ratio to the baseline spans
from $\times 1.30$ to $\times 2.42$.

\begin{keybox}[title={The practical rule}]
Report the baseline in the same table, under the same protocol. It does not
repair the protocol, and it does make your number readable by someone whose
protocol differs --- which is everyone.
\end{keybox}

\subsection{What the numbers actually license}

Under the honest protocol --- temporal, full catalogue --- the factorisation
scores 0.0926 and popularity 0.0404. The model is worth 2.3 times the baseline,
and \emph{that} is the real result. Every other number in the table is larger
and means less. NDCG@10 tells the same story: 0.0464 against 0.0196.

\begin{keybox}[title={The honest headline}]
``A rank-32 factorisation places the user's next film in the top ten 9.3\% of
the time, against 4.0\% for a popularity list, under a temporal split with the
full catalogue ranked.''

Every clause is load bearing.
\end{keybox}

\subsubsection{Why the honest numbers look so small}

9.3\% sounds like failure. It is not: the task is to place one specific film in
the top ten of 3{,}706, from a history that ends before it.

\begin{center}
\small
\begin{tabular}{lr}
\toprule
Random ranker  & 0.0027 \\
Popularity     & 0.0404 \\
Factorisation  & 0.0926 \\
\bottomrule
\end{tabular}
\end{center}

Against chance the factorisation is 34 times better. Small absolute numbers are
what an honest protocol on a hard task looks like --- and a field that finds
small numbers embarrassing will drift toward protocols that produce large ones.
That drift is not fraud, and it does the same damage.

\subsubsection{The mean hides the users}

HR@10 per user is 0 or 1, so a mean of 0.0926 means precisely this: the
held-out film was in the top ten for 9.3\% of users and not for the rest. There
is no partial credit and no distribution to inspect --- the mean \emph{is} the
proportion, which makes it honest and makes its confidence interval computable,
since it is a proportion over 6{,}040 users. NDCG@10 does have a spread, and it
is rarely reported. ``Report the variance'' is good advice that means different
things for different metrics: for HR the right companion is an interval, for
NDCG a distribution.

\subsection{So what should you actually do?}

\begin{enumerate}[leftmargin=*,nosep]
  \item Rank the full catalogue. It is affordable now; the shortcut is a
        historical artefact.
  \item Split in time, per user, and say so.
  \item Include popularity and a tuned classical baseline in the same table.
  \item Never quote a number from another paper into your table without
        re-running it under your protocol.
  \item Report coverage and the popularity profile beside the accuracy metric.
\end{enumerate}

All five are cheap. None of them is standard. Every one of them can only make
your headline number smaller, which is the whole difficulty.

\section{Popularity bias, and what an offline number is worth}

\subsection{The feedback loop}

A recommender is not an observer of its data. It \emph{produces} it. The system
recommends popular items because they are the safe bet; users interact with
what they are shown, since they cannot interact with what they are not; those
interactions become the next training set; and the popular items become more
popular still.

\begin{failbox}[title={Nothing in the objective notices}]
Every offline metric in this lecture is computed on interactions produced by
whatever system was running at the time. A model that reproduces the old
system's behaviour scores well, and a model that would have shown something
better scores badly --- because the evidence for that something was never
collected.
\end{failbox}

\subsubsection{It is the pooling problem again}

\begin{center}
\small
\begin{tabular}{p{0.42\textwidth}p{0.42\textwidth}}
\toprule
\textbf{\lecture{19}} & \textbf{\lecture{22}} \\
\midrule
unjudged documents counted as irrelevant & unseen items counted as negatives \\
the pool came from the systems of its day & the log came from the recommender
of its day \\
a novel method is penalised for finding new documents & a novel model is
penalised for recommending new items \\
\bottomrule
\end{tabular}
\end{center}

Three lectures apart, the same structure: the ground truth records what a
previous system chose to show, and everything else is scored as wrong. Which is
why Part~V ends here rather than on an architecture. The architectures
converged two lectures ago; the measurement problem did not.

\subsection{What to report, beside the metric}

\begin{itemize}[leftmargin=*,nosep]
  \item \term{Catalogue coverage} --- what fraction of items are ever in
        anyone's top ten. A model with a great HR@10 and 2\% coverage is a
        popularity list with extra steps.
  \item \term{Popularity profile} --- the distribution of recommended items'
        interaction counts, against the catalogue's.
  \item \term{Per-user spread}, not just the mean.
  \item \term{The protocol}, in the sentence that states the number.
\end{itemize}

None of these is hard. They are absent from papers because nothing forces them
to be present, and because each of them can only make the headline look worse.

\subsection{None of it is the real test}

Every number in this lecture is \emph{offline}: computed on a log, about a
system that was never run.

\begin{center}
\small
\begin{tabular}{p{0.42\textwidth}p{0.42\textwidth}}
\toprule
\textbf{Offline} & \textbf{Online} \\
\midrule
did the model rank the held-out item highly? & did people click, watch, return,
cancel? \\
on a log the old system produced & on traffic the new system produces \\
a proxy, computable in seconds & the objective, over weeks \\
no risk & real users see the worse arm \\
\bottomrule
\end{tabular}
\end{center}

The last row is why offline evaluation exists at all, and the second is why it
cannot settle anything: the offline test asks whether you would have reproduced
the old system's successes.

\begin{keybox}[title={The honest position}]
Use offline evaluation to decide what is worth testing. Use an online test to
decide what is worth keeping. Do not use either to make the claim the other one
supports.
\end{keybox}

\subsection{One more thing the metric cannot see}

A recommender chooses what a person is exposed to, at scale, and the objective
it is optimised for is engagement measured by interaction. Content that
provokes engages, and engagement is the signal the model is fitted to; the
feedback loop then amplifies whatever the objective rewards, including things
nobody chose to reward. None of this appears in HR@10, NDCG@10, coverage, or
any metric in this lecture.

This course cannot settle what a recommender should optimise. It can insist
that you know your objective is a \emph{proxy}, that you can say what it omits,
and that ``the metric improved'' is never by itself an argument that the system
got better.

\subsection{Sequential recommendation, briefly}

Everything so far treated a user as a \emph{set} of interactions. But order
carries information: someone who has just watched the first two films of a
trilogy is in a \emph{state}, not merely in possession of a history. Model the
sequence and predict the next item --- which is the language-modelling
objective of \lecture{17} with items as tokens --- using architectures you
already have: a recurrent network (\lecture{16}) or self-attention over the
history (\lecture{18}). And a temporal split becomes \emph{mandatory}, because
a random one destroys the very ordering the model is about.

\subsection{The convergence, in one table}

\begin{center}
\small
\begin{tabular}{p{0.16\textwidth}p{0.32\textwidth}p{0.34\textwidth}}
\toprule
 & \textbf{Search (\lecture{20})} & \textbf{Recommendation (\lecture{22})} \\
\midrule
Left tower   & query & user \\
Right tower  & document & item \\
Score        & dot product & dot product \\
Negatives    & in-batch, hard-mined & in-batch, sampled \\
Index        & ANN over documents & ANN over items \\
Ground truth & pooled judgements & a log from the old system \\
\bottomrule
\end{tabular}
\end{center}

Six rows, and only the last two differ in substance. The two halves of Part~V
are one method with two data-collection problems, and the data-collection
problems are the hard part in both.

\begin{trybox}[title={Exercise 22.1}]
Drop the sampled negatives from 100 to 20 and re-run. Predict the direction
first, then check whether the \emph{ordering} of the two methods survives.
\end{trybox}

\section{Reference}

\subsection{Five questions for a recommender paper}

\begin{enumerate}[leftmargin=*]
  \item \term{Full catalogue or sampled?} If sampled, the number is not
        comparable to anything.
  \item \term{Random split or temporal?} A random split leaks the user's
        future.
  \item \term{Is popularity in the table}, under the same protocol?
  \item \term{Were the baselines tuned} as carefully as the proposed model, or
        quoted from elsewhere?
  \item \term{Is there an online result?} If not, the claim is about a log, not
        about users.
\end{enumerate}

In the reproducibility study, most papers failed at least two.

\subsubsection{An exam-style question, worked}

\emph{``A paper reports HR@10 of 0.8085 for its new model, against 0.0404 for a
popularity baseline quoted from an earlier paper. What, if anything, has been
established?''}

\begin{itemize}[leftmargin=*,nosep]
  \item Nothing. The two numbers are from different protocols, and a quoted
        baseline is by definition not run under yours.
  \item Under one protocol, popularity itself reaches 0.6201 --- larger than
        the honest counterpart of the number being celebrated.
  \item The repair costs nothing: re-run popularity under the paper's own
        protocol and put it in the same table.
\end{itemize}

\subsection{Where students lose marks}

\begin{itemize}[leftmargin=*,nosep]
  \item Comparing HR@10 figures from two papers without checking the protocol.
        This lecture exists to make that reflex automatic.
  \item Saying sampled metrics are ``noisier but unbiased''. They are biased,
        and non-uniformly across models.
  \item Omitting the popularity baseline.
  \item Forgetting to remove already-seen items before ranking.
  \item Describing a two-tower model as fundamentally different from matrix
        factorisation. It is the same score with richer encoders.
\end{itemize}

\subsection{Vocabulary}

\begin{center}
\small
\begin{tabular}{ll}
\toprule
Two-tower        & user encoder and item encoder, scored by a dot product \\
Sampled softmax  & a softmax whose denominator is estimated from a sample \\
$\log q$ correction & subtracting $\log q(j)$ to unbias that estimate \\
Leave-one-out    & hold out one interaction per user, chosen at random \\
Temporal split   & hold out the most recent interaction instead \\
Sampled metric   & rank the held-out item against a sample of negatives \\
Popularity bias  & systematic over-recommendation of already-popular items \\
Feedback loop    & the system shaping the data it will next be trained on \\
\bottomrule
\end{tabular}
\end{center}

\subsection{Scope}

\term{Examinable.} Why RMSE is the wrong metric for top-$N$; the two-tower
model and its identity with \lecture{20}'s bi-encoder; sampled softmax and the
$\log q$ correction; the four protocols, the direction of each effect, and why
sampling can reverse an ordering; the popularity baseline; popularity bias and
the feedback loop; the offline/online distinction.

\term{Not examinable --- engineering.} Specific architectures and libraries,
serving, A/B testing machinery.

\term{Beyond the syllabus.} Sequential and session-based models in detail,
counterfactual and off-policy evaluation, bandits for exploration.

\subsection{Reading}

\begin{itemize}[leftmargin=*,nosep]
  \item These notes are the primary source.
  \item Dacrema, Cremonesi and Jannach, \emph{Are We Really Making Much
        Progress?} --- the reproducibility study. Read it before you believe
        any recommender leaderboard.
  \item Krichene and Rendle, \emph{On Sampled Metrics for Item Recommendation}
        --- why sampled metrics are not merely noisy.
  \item Yi et al., \emph{Sampling-Bias-Corrected Neural Modeling} --- the
        $\log q$ correction in a deployed two-tower system.
\end{itemize}

\subsection{Three numbers to leave with}

\begin{center}
\small
\begin{tabular}{lr}
\toprule
Popularity, honestly measured                & 0.0404 \\
A rank-32 factorisation, honestly measured   & 0.0926 \\
Popularity, measured generously              & 0.6201 \\
\bottomrule
\end{tabular}
\end{center}

Rows one and two are the finding: a real model is worth 2.3 times the trivial
baseline. Row three is that \emph{same trivial baseline}, scored under a
protocol many papers use, reporting 6.7 times the honest model. If you remember
one thing from Part~V, make it that row three exists and looks exactly like a
result.

\subsection{Part V, closed}

\begin{center}
\small
\begin{tabular}{p{0.16\textwidth}p{0.68\textwidth}}
\toprule
\lecture{19} & a ranking is scored by a question, chosen before the scores \\
\lecture{20} & the first stage sets the ceiling; the old method still won \\
\lecture{21} & the missing entries are the data; do not impute them \\
\lecture{22} & the protocol is part of the result \\
\bottomrule
\end{tabular}
\end{center}

Four lectures outside the textbook, and not one of them needed a method you had
not already met. What they needed was the discipline of Part~I applied to
problems whose ground truth is itself a measurement.

\begin{keybox}[title={The Part V checklist}]
\begin{enumerate}[leftmargin=*,nosep]
  \item Name the metric and the cutoff before seeing scores.
  \item Fit the trivial baseline first --- corpus order, popularity, the global
        mean.
  \item State the protocol in the same sentence as the number.
  \item Ask where the ground truth came from, and what it counts as negative.
  \item Report what the metric cannot see, in one line.
\end{enumerate}
\end{keybox}

\subsection{Summary}

\begin{keybox}[title={One line to keep}]
The protocol is part of the result. A number without it is not a measurement.
\end{keybox}

A recommender ranks a catalogue for a user, so it is scored by \lecture{19}'s
metrics on items the user had not touched --- not by RMSE. The two-tower model
is \lecture{20}'s bi-encoder with users in place of queries, and matrix
factorisation is its lookup-table special case. The catalogue softmax is
unaffordable, so it is sampled --- and a sampled softmax needs the $\log q$
correction or popularity bias enters the gradient. The protocol moved HR@10
from 0.0926 to 0.8085 on one unchanged model, and it lets a popularity list
report 6.7 times a real model's honest score. And offline numbers filter ideas;
only an online test measures the thing you care about.

\end{document}
